\section{Experimental Methodology}
\label{sec:experimental-methodology}

\subsection{Environment, metrics, and statistics}

The retained experiments ran in the project's Docker/WSL environment using project container images, Python runners, Linux traffic control, ping, and iperf3. Exact host hardware utilization was not retained consistently, so hardware-level performance generalization is avoided. Every reported result is tied to a run directory or repository summary in the evidence inventory.

The primary metrics are TCP throughput in Mbps, ping RTT in milliseconds, end-to-end ping packet loss in percent, and controller reward. For complete per-run cohorts, the audit uses the sample standard deviation and a two-sided Student's $t$ 95\% confidence interval,
\[
  \bar{x} \pm t_{0.975,n-1}\frac{s}{\sqrt{n}}.
\]
Values already present in repository summaries are labelled repository-reported (RR); statistics recomputed without changing the evidence repository are labelled audit-computed (AC). No missing measurements are imputed.

\begin{table}[h]
\caption{RQ1 analysis matrix. A configured token rate is a reference target, not an experimental baseline.}
\label{tab:experimental-matrix-rq1}
\centering
\footnotesize
\setlength{\tabcolsep}{1pt}
\begin{tabular}{p{0.12\linewidth}p{0.21\linewidth}p{0.20\linewidth}p{0.16\linewidth}p{0.22\linewidth}}
\hline
Case & Control/reference and variable & Outcome & Valid/excluded & Analysis/evidence \\
\hline
BW20 & No separate baseline; configured token rate is target & TCP rate & 5/0 & RR mean, SD, CI, range; balanced bandwidth \\
BW50 & No separate baseline; configured token rate is target & TCP rate & 5/0 & RR mean, SD, CI, range; balanced bandwidth \\
BW100 & No separate baseline; configured token rate is target & TCP rate & 5/0 & RR mean, SD, CI, range; balanced bandwidth \\
Delay & 0-ms control; 0/10/30/50-ms one-way delay & RTT & 5 each/0 & Overall permutation test, epsilon-squared, AC intervals; latency \\
Loss & 0\% control; 0/1/3/5\% one-way loss; model round-trip reference & End-to-end loss & 5 each/0 & Overall permutation test, epsilon-squared, AC intervals/counts; loss \\
\hline
\end{tabular}
\end{table}

\begin{table}[h]
\caption{RQ2/RQ3 and interface analysis matrix. These cases remain descriptive.}
\label{tab:experimental-matrix-rq23}
\centering
\footnotesize
\setlength{\tabcolsep}{1pt}
\begin{tabular}{p{0.12\linewidth}p{0.23\linewidth}p{0.20\linewidth}p{0.16\linewidth}p{0.20\linewidth}}
\hline
Case & Control/reference and variable & Outcome & Valid/excluded & Analysis/evidence \\
\hline
Dual router & Zero impairment; fixed topology & TCP rate & 5/0 & RR/AC descriptive; dual-router \\
Germany paths & Representative path case; hops and cumulative delay co-vary & RTT; TCP rate & 3 each/0 & Confounded case comparison; Germany50 paths \\
Germany full & Selected execution vs full-plan dry-run & Route/traffic counts & 9 traffic + 1 dry-run/0 & Counts/pass-fail; Germany50 full \\
AI & Provider/validation path; no performance baseline & Validation and one-run metrics & 1 + 1/0 & Categorical/descriptive; official/Qwen \\
Control & Fixed A/B and heuristic comparators; policy & Reward and network metrics & 80/2 retries & Single-seed ordered measurement rows; RR descriptive summary; RL \\
Dashboard & No baseline; one direct scenario & HTTP, files, metrics & 1 documented/0 & Documented-only descriptive record \\
\hline
\end{tabular}
\end{table}


\subsection{Controlled impairment experiments}

The bandwidth evidence has deliberately unequal provenance. The direct 20-Mbps result is one retained historical client--router--server measurement using server-side TBF. It is reported as $n=1$. The 50- and 100-Mbps results are later supplementary replications of the direct two-node, server-side TBF, five-second reverse-TCP iperf3 method. The 50-Mbps cohort contains five valid runs. The 100-Mbps cohort contains five valid measurements from six attempts; the failed attempt and explicit retry remain in the evidence.

Latency used a 20-Mbps routed topology at configured one-way delays of 0, 10, 30, and 50 ms. Each condition has five valid runs and three ping packets per run. The dependent variable is measured RTT, not one-way latency. Packet-loss experiments configured one-way loss of 0, 1, 3, and 5\%, with five runs and 100 ping packets per run. Sent and received counts determine measured end-to-end loss. The dual-router benchmark used client--router1--router2--server, zero configured delay/loss, and five runs; it is analyzed separately from direct bandwidth evidence.

\subsection{Topology, AI, control, and interface experiments}

Three Germany50-derived paths were fixed before execution: Aachen--Koeln (one backbone hop), Erfurt--Wuerzburg--Augsburg--Muenchen--Passau (four), and Oldenburg--Bremen--Hannover--Braunschweig--Kassel--Erfurt--Wuerzburg--Augsburg--Muenchen--Passau (nine). Extracted-path experiments used 20 Mbps, 10-ms one-way delay, and 0\% loss on each link, with three runs per path. A separate full-topology evaluation instantiated all 50 nodes and 88 links. Selected routing installed 20 routes for three endpoint pairs and produced nine real traffic measurements. Full routing generated and dry-run validated 4,224 entries but did not install or exercise them as all-pairs traffic.

AI validation distinguishes the official OpenAI attempt from a separate OpenAI-compatible Qwen path. The official request produced HTTP 429 \texttt{insufficient\_quota}. The Qwen-compatible path produced a six-node scenario that passed schema, semantic, prompt-constraint, projection, dry-run, and controlled real-Docker validation.

The RL comparison contains 80 valid real-Docker episodes: 20 each for Q-learning, threshold heuristic, fixed A, and fixed B. All policies experienced the same ordered two-phase schedule. The recorded Q-learning configuration used seed 20260722, learning rate $\alpha=0.35$, discount $\gamma=0.85$, and $\epsilon=0.18$. Two failed trials were retained as failed/retried checkpoints and are not counted among the 80 valid rows. The Dashboard result is a documented single direct UI-path smoke, not a repeated experiment; absence of its raw validation logs determines its documented-only status.
